\documentclass[../../main.tex]{subfiles}
\begin{document}

\begin{center}
  {\fontsize{48pt}{58pt}\selectfont \textbf{1992}\normalsize\textbf{年度}}
\end{center}

\begin{table}[htb]
  \centering
  \caption{出題テーマと難易度}
  \begin{tabular}{|c|p{9cm}|c|}
    \hline
    \textbf{問題番号} & \textbf{テーマ} & \textbf{難易度} \\
    \hline
    第1問 & 対数曲線に内接する四辺形の面積 & \\
    \hline
    第2問 & 格子点を頂点にもつ三角形 & \\
    \hline
    第3問 & 四面体の内接球の半径と不等式 & \\
    \hline
    第4問 & 円柱面を切り開いた展開図 & \\
    \hline
    第5問 & 曲線上の運動と速度ベクトル & \\
    \hline
    第6問 & じゃんけんの期待値と戦略 & \\
    \hline
  \end{tabular}
\end{table}

\begin{table}[htb]
  \centering
  \caption{解答}
  \renewcommand{\arraystretch}{1.4}
  \begin{tabular}{|c|c|p{9cm}|}
    \hline
    \textbf{問題番号} & \textbf{小問} & \textbf{答え} \\
    \hline
    第1問 & (1) & $b=\dfrac{a-1}{\log a}$，$S_2=\dfrac{a\log a-a+1}{2}+\dfrac{a-1}{2}\log\left(\dfrac{a-1}{\log a}\right)$ \\
    \cline{2-3}
    & (2) & 極限値は$1$ \\
    \hline
    第2問 & (1) & 辺$BC$の内部にも奇数個の格子点がある（証明問題） \\
    \cline{2-3}
    & (2) & 面積は$8$の倍数（証明問題） \\
    \hline
    第3問 & -- & 2つの不等式が成立（証明問題） \\
    \hline
    第4問 & -- & 展開図は解答の図を参照 \\
    \hline
    第5問 & (1) & $t^2=2\left(\dfrac{x^3}{3}-\dfrac{1}{x}+\dfrac{2}{3}\right)$ \\
    \cline{2-3}
    & (2) & $\sqrt{\dfrac{17}{3}}$秒後 \\
    \hline
    第6問 & (1) & チョキを確率$1$で出す \\
    \cline{2-3}
    & (2) & グー，チョキ，パーを出す確率は$\left(\dfrac{5}{14},\dfrac{3}{7},\dfrac{3}{14}\right)$ \\
    \hline
  \end{tabular}
\end{table}

\paragraph{試験形態}
6問．試験時間は要確認．

\end{document}
